% page 186 of the facsimile (image p-185 of the scan)
% block 157.5 x 242.0 mm ; left margin 8.0 mm
\pgc{-12.700mm}{0.000mm}{
\l{\cel{4}178}
\l{}
\l{}
\l{\sou{francisk}o.(\sou{arkeol.)}\cel{2}Hakilo militala, kun du tranchivi,}
\l{\cel{3}uzita da la Franki antiqua. - DEFS.}
\l{}
\l{\sou{franciskano}. Reguliero di la ordeno (religiala, katolika)}
\l{\cel{3}quan fondis santa Franciscus di Assisi. - DEFIRS.}
\l{}
\l{\sou{frangulo.}\cel{2}(\sou{bot.}) Arboreto qua kreskas inter la bushi e la}
\l{\cel{3}hegi, en la tereni humida, di qua la ligno esas transfor-}
\l{\cel{3}mebla a karbono lejera quan onu uzas por fabrikar la kanon-}
\l{\cel{3}pulvero, di qua la kortico interna esas purgiva, e di qua}
\l{\cel{3}la beri esas verda o reda, segun la grado di matureso.}
\l{\cel{3}- L.\cel{1}\sou{rhamnus frangula}. - IL.}
\l{}
\l{\sou{franjo.}\cel{2}I. Ornivo ek serio de brini, de torsadi pendanta}
\l{\cel{3}(kotona, lana, silka, orfila, arjenta, e c.). - II.}
\l{\cel{3}(\sou{fiziko}).Franji, linei alterne briloza ed obskura, en la}
\l{\cel{3}fenomeno \textquotedbl{}interfero\textquotedbl{}. - DEFIS.}
\l{}
\l{\sou{franjipano}.\cel{2}Sorto di aromo. - DEFIS.}
\l{}
\l{\sou{franko}. Unajo monetala di la sistemo decimala qua korespondas}
\l{\cel{3}a moneto-peco ek arjento aloyita a 0 fr.165 de kupro, qua}
\l{\cel{3}pezas 5 grami, dividita a 10 dekimi o 100 centimi. - F.}
\l{}
\l{\sou{frankolino}. (\sou{zool.)}\cel{3}Ucelo de la genero \textquotedbl{}perdriki\textquotedbl{}, sama-}
\l{\cel{3}statura kam la fazano, di qua la gambi esas longa, la beko}
\l{\cel{3}forta e longa. - L.\cel{1}\sou{pternistes vulgaris}.\cel{2}DEFIRS.}
\l{}
\l{\sou{frapar}.\cel{2}(\sou{trans.}) Shokar unfoye o plurafoye per apliko vio-}
\l{\cel{3}lentoza. - F.}
\l{}
\l{\sou{frato}.\cel{2}I. Persono konsiderata relate lua parenteso kun altru}
\l{\cel{3}qua naskis de la sama gepatri. - II. (\sou{metaf}.)Persono konsi-}
\l{\cel{3}derata relate la ligilo qua unionas lu a la membri cetera}
\l{\cel{3}di la familio homara. - EFIS.}
\l{}
\l{\sou{fraudar}. (\sou{trans.)}\cel{2}I. Trompar (ulu) por prokurar a su, detri-}
\l{\cel{3}mento di ta ulu, ca o ta avantajo. - II. Frustrar la admi-}
\l{\cel{3}nistrantaro doganala o fiskala. - DEFIS.}
\l{}
\l{\sou{fraxino}. (\sou{bot.}) Arboro de la familio \textquotedbl{}oleacei\textquotedbl{}, di qua la}
\l{\cel{3}ligno, kompakta, blanka e veinoza, uzesas en industrio.}
\l{\cel{3}L.\cel{1}\sou{fraxinus}. FISL.}
\l{}
\l{\sou{frayo}. Fish-ovi fekundigita. - FIS.}
\l{}
\l{\sou{frazo}. I. (\sou{gram.}) Propoziciono simpla od asemblajo de propo-}
\l{\cel{3}zicioni qua facas senco kompleta. - II. (\sou{muziko}). Serio de}
\l{\cel{3}son-uni qui prizentas a la oreli developeso reguloza e}
\l{\cel{3}kompleta. - DEFIRS.}
\l{}
\l{\sou{frazeologio}. Maniero konstruktar sua frazi, qua karakterizas}
\l{\cel{3}linguo o skriptisto. - DEFIRS.}
}
